\chapter{Index Construction and Rebalancing}\label{ch:m1:index-construction-and-rebalancing}

At the close of 30 April 2026 nothing visible happened: the capitalisation of
every listed American company was written down. Eight weeks later, after the
close of 26 June, the funds that follow one family of indices all traded to
the list that snapshot had produced, in the same stocks, in the same
direction, in the same auction. Everyone had known the rule for a year and the
list for a month. An index is a published algorithm with a calendar, and a
very large amount of money has promised to obey it. This chapter reads the
algorithm as a trader does: what will have to be bought, by whom, when, and
how much of the consequence is already in the price.

\section{An index is a rulebook}

\begin{definition}[Index methodology]\label{def:m1:index-construction-and-rebalancing:methodology}
An \emph{index methodology}\index{index methodology} is the published
document that determines an index: the universe of eligible securities, the
selection rule, the weighting rule, the calendar of reviews and the treatment
of corporate actions. Selection is \emph{rules-based} when the document
alone determines the members, and \emph{committee-based} when a committee
chooses among the eligible.
\end{definition}

\begin{definition}[Float adjustment]\label{def:m1:index-construction-and-rebalancing:floatadj}
\emph{Float adjustment}\index{float adjustment} weights each member by its
capitalisation multiplied by a float factor, the fraction of its shares in
the \omterm{def:m1:shares-corporate-actions-indices:float}{free float} (\cref{def:m1:shares-corporate-actions-indices:float}), so
that the index can be held by everyone at once.
\end{definition}

\begin{definition}[Reconstitution and rebalance]\label{def:m1:index-construction-and-rebalancing:recon}
An \emph{index reconstitution}\index{index reconstitution} is a scheduled
review at which the membership is redetermined. A \emph{rebalance} updates
the weights (share counts, float factors, caps) of an unchanged membership.
Both follow a fixed sequence: a \emph{rank day} on which the data are
frozen, an \emph{announcement}, and an \emph{effective date}, after whose
close the new index applies.
\end{definition}

\begin{definition}[Pro-forma index]\label{def:m1:index-construction-and-rebalancing:proforma}
A \emph{pro-forma index}\index{pro-forma index} is the index as it will
stand after a coming review, members and share counts, distributed by the
provider to its clients between the announcement and the effective date so
that trackers can prepare their trades.
\end{definition}

\begin{dated}{2026-09}{Two ways of choosing five hundred or two thousand stocks}\label{dat:m1:index-construction-and-rebalancing:rules}
\textbf{Rules-based.} The Russell US indices rank companies by total
capitalisation on a rank day. From 2026 they reconstitute twice a year. June
2026: rank day 30 April, preliminary lists updated on 29 May and 5, 12 and 18
June, effective after the close of Friday 26 June; the breakpoint between the
large-capitalisation Russell 1000 and the small-capitalisation Russell 2000
was \$5.7 billion, 24\% higher than a year before. December 2026, the first
December reconstitution in more than three decades: rank day 30 October,
preliminary lists 13 November, effective after the close of 11 December.
The provider puts the assets benchmarked to or invested in products based on
these indices at about \$12.2 trillion.

\textbf{Committee-based.} The S\&P 500's members are chosen by a committee
among companies that meet published criteria. Since 1 July 2025 an addition
needs a total capitalisation of at least \$22.7 billion (the guideline is
reviewed each quarter and set near the 85th percentile of cumulative \omterm{def:m1:shares-corporate-actions-indices:share}{market
capitalisation}), a float-adjusted capitalisation of at least half that
threshold, and a ratio of annual value traded to float-adjusted
capitalisation of at least 0.75, besides tests of profitability and
domicile. The criteria are for additions, not for continued membership.
\end{dated}

\begin{omfigure}
\begin{tikzpicture}[font=\small]
\draw[thick,-{Stealth[length=2.5mm]}] (0,0) -- (13,0);
\foreach \x/\d/\t in {0.8/30 April/rank day,5.0/29 May/updated list,10.9/26 June/effective after the close}{
  \draw[thick] (\x,0.12) -- (\x,-0.12);
  \node[above=3pt,font=\footnotesize\bfseries] at (\x,0.12) {\d};
  \node[below=3pt,font=\footnotesize,align=center] at (\x,-0.12) {\t};}
\foreach \x in {6.2,7.4,8.6,9.6}{\draw (\x,0.08) -- (\x,-0.08);}
\node[above=16pt,font=\footnotesize] at (7.9,0.12) {further updates: 5, 12, 18 June};
\draw[decorate,decoration={brace,amplitude=5pt,mirror},omDef,thick] (0.8,-1.0) -- (4.9,-1.0) node[midway,below=6pt,font=\footnotesize,align=center,text=omDef] {anyone can compute the list:\\prediction};
\draw[decorate,decoration={brace,amplitude=5pt,mirror},omThm,thick] (5.1,-1.0) -- (10.9,-1.0) node[midway,below=6pt,font=\footnotesize,align=center,text=omThm] {the list is public:\\positioning, then the closing auction};
\end{tikzpicture}

\omcaption{The June 2026 reconstitution of a rules-based index family
(dates from the dated box). Prices after the rank day do not change the
membership; they change only how much the trackers will have to
buy.}\label{fig:m1:index-construction-and-rebalancing:timeline}
\end{omfigure}

\section{Buffers}

A pure ``top $n$ by capitalisation'' rule makes the stocks near rank $n$ enter
and leave at every review. Each change costs the trackers a trade, so
providers damp the rule.

\begin{definition}[Buffer rule]\label{def:m1:index-construction-and-rebalancing:buffer}
A \emph{buffer rule}\index{buffer rule} gives incumbents an advantage at the
boundary: with a buffer $b$, a current member of a top-$n$ index remains while
it ranks $n+b$ or better, and a non-member enters only when it ranks $n-b$ or
better; any places left are filled by rank. Providers express $b$ in ranks,
in percentiles of cumulative capitalisation, or as a percentage band around
the breakpoint.
\end{definition}

\begin{omfigure}
\begin{tikzpicture}
\begin{axis}[width=11cm,height=6.2cm,
  xlabel={rank at this review},ylabel={rank at the last review},
  tick label style={font=\small},label style={font=\small},
  xmin=148,xmax=252,ymin=80,ymax=420,grid=major,grid style={gray!25},
  legend columns=3,legend cell align=left,legend style={font=\footnotesize,at={(0.5,-0.27)},anchor=north,draw=none,fill=none,/tikz/every even column/.append style={column sep=6pt}},
  scatter/classes={stays={mark=*,omExo!70,mark size=1.4pt},outside={mark=o,omExo!70,mark size=1.4pt},
    added={mark=triangle*,omDef,mark size=2.6pt},deleted={mark=square*,omThm,mark size=2pt},
    kept={mark=*,omThm,mark size=2pt},blocked={mark=o,omDef,thick,mark size=2pt}},
  table/col sep=comma]
\addplot[scatter,only marks,scatter src=explicit symbolic] table[x=new_rank,y=old_rank,meta=outcome]{figdata/markets-1/15-index-construction-and-rebalancing/band.csv};
\draw[thick,dashed] (axis cs:200,80) -- (axis cs:200,420);
\draw[thin] (axis cs:180,80) -- (axis cs:180,420);
\draw[thin] (axis cs:220,80) -- (axis cs:220,420);
\legend{member stays,outside,added,deleted,member kept by the buffer,outsider blocked by the buffer}
\end{axis}
\end{tikzpicture}

\omcaption{One simulated review of a top-200 index with a buffer of 20
ranks (thin lines), for the stocks now ranked 150 to 250. Between ranks 180
and 220 history decides: members stay, outsiders wait. Data: the tutorial's
simulation.}\label{fig:m1:index-construction-and-rebalancing:band}
\end{omfigure}

\begin{omfigure}
\begin{tikzpicture}
\begin{axis}[width=9.5cm,height=4.6cm,
  xlabel={buffer (ranks)},ylabel={additions per review},
  tick label style={font=\small},label style={font=\small},
  xmin=-3,xmax=63,ymin=0,ymax=24,xtick={0,10,20,30,40,60},grid=major,grid style={gray!25},
  nodes near coords,nodes near coords style={font=\footnotesize,above=2pt},
  point meta=explicit,table/col sep=comma]
\addplot[omDef,very thick,mark=*] table[x=buffer,y=names,meta expr=\thisrow{weight_pct}]{figdata/markets-1/15-index-construction-and-rebalancing/turnover.csv};
\end{axis}
\end{tikzpicture}

\omcaption{Average number of additions per annual review of a simulated
top-200 index from a universe of 600, against the width of the buffer. The
number above each point is the one-way turnover in percent of index weight. A
buffer of 30 ranks halves the number of changes. Data: the tutorial's
simulation, 20 paths of 20 years.}\label{fig:m1:index-construction-and-rebalancing:turnover}
\end{omfigure}

\begin{method}[Predicting a rules-based reconstitution]\label{met:m1:index-construction-and-rebalancing:predict}
\begin{enumerate}
\item \textbf{Universe.} Rebuild the eligible universe from the methodology:
listing venue, domicile, minimum price and float, share classes.
\item \textbf{Capitalisation.} Compute total capitalisation as the provider
does, with all share classes and the latest share counts from filings. Most
prediction errors are share-count errors.
\item \textbf{Rank and band.} Rank, locate the breakpoints, apply the buffer
to current members.
\item \textbf{Before the rank day}, give each boundary stock a probability
from its distance to the breakpoint and its volatility over the days that
remain.
\item \textbf{Size the trade} with \cref{prop:m1:index-construction-and-rebalancing:demand}.
\end{enumerate}
\end{method}

\section{What the trackers must do}

\begin{proposition}[Passive demand]\label{prop:m1:index-construction-and-rebalancing:demand}
Let funds with assets $A$ replicate a float-adjusted index of total
float-adjusted capitalisation $M$. They hold the same fraction $\varphi =
A/M$ of the float of every member. When a stock with float-adjusted
capitalisation $F$ is added they must buy $\varphi F$ dollars of it, that is
a fraction $\varphi$ of its float (more precisely $AF/(M+F)$, the difference
being negligible for one addition).
\end{proposition}
\begin{proof}
A tracker holds each member at its index weight $F_i/M$, hence $AF_i/M$
dollars of member $i$, which is the fraction $A/M$ of $F_i$.
\end{proof}

\begin{definition}[Demand shock]\label{def:m1:index-construction-and-rebalancing:shock}
A \emph{demand shock}\index{demand shock} is a change in the quantity of a
security that some investors must hold for reasons unrelated to its value. An
index change is the cleanest example: the quantity, the date and the identity
of the buyers are known, and the buyers' benchmark is the closing price of
the effective date, so they are indifferent to the price they pay at that
close.
\end{definition}

\begin{example}[Thirteen days of volume]\label{ex:m1:index-construction-and-rebalancing:thirteen}
An index has $M = \$5$ trillion and trackers with $A = \$2$ trillion, so
$\varphi = 40\%$. A company with 400 million shares at \$50, of which 65\%
float, has $F = \$13$ billion: a weight of 0.26\% and a demand of \$5.2
billion. It trades \$400 million a day: the trackers need thirteen average
days of volume, at one close.
\end{example}

\begin{proposition}[The funding trade]\label{prop:m1:index-construction-and-rebalancing:funding}
When an addition takes weight $w$ in the index, every surviving member's
weight is multiplied by $1-w$: trackers sell $Aw_iw$ of member $i$, and $Aw$
in total less the value of what is deleted.
\end{proposition}

\begin{example}[One large addition]\label{ex:m1:index-construction-and-rebalancing:large}
On 16 November 2020 a provider announced that a company large enough to
weigh more than one percent would join its flagship 500-stock index. It
consulted the market on whether to add it in one step or two, and on 30
November announced a single step, at full float-adjusted weight, effective
before the open of Monday 21 December, the pro-forma files being distributed
after the close of 11 December. With $A = \$10$ trillion and $w = 1.5\%$
(round numbers, not the provider's) the trackers would buy \$150 billion of
one stock and sell \$9 billion of a member weighing 6\%, all at one close:
that of Friday 18 December, the third Friday of the month and so a quarterly
expiry of index futures and options, among the most liquid closes of the
year (\cref{ch:m1:auctions}).
\end{example}

\section{The index effect, and where it went}

\begin{definition}[Index effect]\label{def:m1:index-construction-and-rebalancing:effect}
The \emph{index effect}\index{index effect} is the abnormal return of a
stock between the announcement of its addition to (or deletion from) an index
and the effective date, together with its partial reversal afterwards.
\end{definition}

If investors held a flat demand curve for each stock, a \omterm{def:m1:index-construction-and-rebalancing:shock}{demand shock} that
carries no information would not move the price. For decades it did: a study
of additions to the S\&P 500 finds an average abnormal return of 7.4\% in the
1990s. The same study finds 0.3\% for 2010--2020, and for deletions a large
negative return in the 1990s against 0.1\% in the later decade, although the
assets tracking the index had multiplied in between. The title of the study
is its conclusion: the effect has disappeared from the announcement window.

Three explanations, not exclusive. \emph{Prediction}: additions are now
anticipated, so the price adjusts before the announcement, where an event
study does not look. \emph{Migration}: many additions to the large index
come from the provider's mid-capitalisation index, whose trackers sell what
the large index's trackers buy. \emph{Liquidity supply}: a population of
firms now exists to buy ahead and sell to the trackers at the close. The
\omterm{def:m1:index-construction-and-rebalancing:shock}{demand shock} has not disappeared; its price has been competed down, and
moved earlier.

\begin{omfigure}
\begin{tikzpicture}
\begin{axis}[width=11cm,height=5cm,
  xlabel={trading days from the effective date},ylabel={abnormal return (\%)},
  tick label style={font=\small},label style={font=\small},
  xmin=-15,xmax=30,ymin=-0.5,ymax=9,grid=major,grid style={gray!25},
  legend columns=2,legend style={font=\footnotesize,at={(0.5,-0.3)},anchor=north,draw=none,fill=none,/tikz/every even column/.append style={column sep=8pt}},
  table/col sep=comma]
\addplot[omThm,very thick] table[x=day,y=then_pct]{figdata/markets-1/15-index-construction-and-rebalancing/event.csv};
\addplot[omDef,very thick,dashed] table[x=day,y=now_pct]{figdata/markets-1/15-index-construction-and-rebalancing/event.csv};
\legend{scaled to the 1990s average (7.4\%),scaled to the 2010--2020 average (0.3\%)}
\end{axis}
\end{tikzpicture}

\omcaption{A stylised addition, in cumulative abnormal return: a jump at the announcement (day $-5$), a
run-up to the effective close, a partial reversal. Only the two heights at
day 0 are published averages; the shapes, the split between permanent and
temporary parts and the speed of the reversal are
illustrative.}\label{fig:m1:index-construction-and-rebalancing:event}
\end{omfigure}

Index events remain a business (One Quant Book 8 gives the strategy files)
because the averages hide a wide cross-section: small-capitalisation
reconstitutions with demands of several days' volume, migrations, float and
share-count changes that nobody announces in a headline, and indices in
markets where fewer firms do the arithmetic.

\section{Tutorial: a reconstitution with buffers}\label{tut:m1:index-construction-and-rebalancing:recon}

\begin{tutorial}
\textbf{Goal.} Reconstitute a top-$n$ index with a buffer, measure what the
buffer saves, and size the trackers' demand. \textbf{End state:} the three
data figures of this chapter.
\begin{tutsteps}
\item \textbf{The rule.} Keep members inside $n+b$, admit outsiders inside
$n-b$, then fill or trim to exactly $n$.
\omcode{code/markets-1/15-index-construction-and-rebalancing/python/index_recon.py}{15}{31}{A top-$n$ reconstitution with a buffer of $b$ ranks.}\label{lst:m1:index-construction-and-rebalancing:recon}
\item \textbf{Turnover.} Simulate 600 lognormal capitalisations with 35\%
annual volatility and review once a year
(\cref{fig:m1:index-construction-and-rebalancing:turnover}).
\item \textbf{Demand.}
\omcode{code/markets-1/15-index-construction-and-rebalancing/python/index_recon.py}{53}{59}{What the trackers must buy, in dollars and in days of volume.}\label{lst:m1:index-construction-and-rebalancing:demand}
\end{tutsteps}
\textbf{What to change next.} Replace the rank buffer by a band of $\pm
2.5\%$ of cumulative capitalisation around the breakpoint and compare the
turnover. Then give each boundary stock a probability of crossing between
today and the rank day, from its volatility, and check the calibration of
your probabilities over the simulated years.
\end{tutorial}

\section{Build: the reconstitution predictor}\label{bld:m1:index-construction-and-rebalancing:predictor}

\begin{build}
\bfield{Purpose} The miniature firm's event book (One Quant Book 8) starts
from a list: which names will enter and leave, and how many days of volume
the trackers will need in each.
\bfield{Interface} \texttt{Security(symbol, shares, price, float\_factor,
adv\_shares, member)}; \texttt{predict(universe, n, buffer)} returning
members, additions and deletions; \texttt{weights(universe, members)},
float-adjusted; \texttt{tracker\_trades(universe, prediction,
tracked\_assets)} returning for each addition and deletion the dollars, the
shares and the multiple of average daily volume.
\bfield{Rules} Ranking on total capitalisation, weighting on float-adjusted
capitalisation: the two are different on purpose. Ties are broken by symbol so
that the output is deterministic. Share counts come from the corporate-action
adjuster (\cref{bld:m1:shares-corporate-actions-indices:adjuster}).
\bfield{Acceptance tests} \texttt{code/firm/recon/tests/}: a five-stock
universe by hand, with and without a buffer; float-adjusted weights; the
trades of the trackers in an addition and a deletion.
\bfield{Stretch} Crossing probabilities before the rank day; migrations
between two indices of the same family with different tracked assets.
\end{build}

\begin{omsources}
\item FTSE Russell, \emph{FTSE Russell begins June 2026 semi-annual Russell
US Indexes reconstitution} and \emph{FTSE Russell announces December 2026
Russell US Indexes reconstitution schedule}, press releases, 2026.
\item S\&P Dow Jones Indices, \emph{Update to S\&P Composite 1500 Market Cap
Guidelines}, press releases of 4 January 2023 and 1 July 2025; \emph{S\&P Dow
Jones Indices Announces Implementation of Tesla's Addition to S\&P 500}, 30
November 2020.
\item R.~Greenwood and M.~Sammon, ``The disappearing index effect'',
\emph{Journal of Finance} 80 (2025), 657--698.
\end{omsources}

\section{Exercises}

\begin{exercise}[$\star$]\label{exo:m1:index-construction-and-rebalancing:1}
A company has 400 million shares at \$50; its founder holds 35\% and the rest
floats. Give its float-adjusted capitalisation and its weight in an index
whose float-adjusted capitalisation is \$5 trillion.
\end{exercise}

\begin{exercise}[$\star$]\label{exo:m1:index-construction-and-rebalancing:2}
Funds with \$2 trillion track that index. What fraction of every member's
float do they hold? The company of the previous exercise is added and trades
\$400 million a day. Give the demand in dollars and in days of volume.
\end{exercise}

\begin{exercise}[$\star$]\label{exo:m1:index-construction-and-rebalancing:3}
A top-100 index has a buffer of 10 ranks. Decide the fate of: a member now
ranked 108; an outsider ranked 95; an outsider ranked 88; a member ranked 112.
\end{exercise}

\begin{exercise}[$\star\star$]\label{exo:m1:index-construction-and-rebalancing:4}
In the June 2026 reconstitution of
\cref{dat:m1:index-construction-and-rebalancing:rules} a company is ranked
just above the \$5.7 billion breakpoint on the rank day and loses 30\% of its
value in May. In which index is it on 29 June? What has the fall changed for
the funds tracking that index, and for a firm that bought the stock in April
expecting their purchases?
\end{exercise}

\begin{exercise}[$\star\star$]\label{exo:m1:index-construction-and-rebalancing:5}
With the round numbers of
\cref{ex:m1:index-construction-and-rebalancing:large}, verify the two dollar
amounts. What is the 6\% member's new weight? Why might the provider have
considered two steps, and what is the argument for one?
\end{exercise}

\begin{exercise}[$\star\star$]\label{exo:m1:index-construction-and-rebalancing:6}
Apply the square-root rule of \cref{met:m1:the-sell-side:sqrt} with
$\kappa = 0.7$ and a daily volatility of 2.5\% to the demand of thirteen
days' volume of \cref{ex:m1:index-construction-and-rebalancing:thirteen}, as
if it were executed in one day. Then suppose anticipating firms buy the same
quantity evenly over the twenty days before and hand it over at the close:
estimate the impact of each of their days. Compare with the two averages
quoted in this chapter.
\end{exercise}

\begin{exercise}[$\star\star\star$]\label{exo:m1:index-construction-and-rebalancing:7}
\emph{Coding.} With \texttt{simulate\_turnover(600, 200, b, 20, 1)} report
the additions per review and the one-way turnover for $b = 0$ and $b = 30$.
An index fund pays 15 basis points on what it trades; what does the buffer
save its investors per year, in basis points of assets?
\end{exercise}

\begin{exercise}[$\star\star\star$]\label{exo:m1:index-construction-and-rebalancing:8}
\emph{Find the flaw.} ``Buying every announced addition to the large index at
the next open and selling at the effective close earned 5\% per event in my
backtest on 1990--2000. There are about twenty-five events a year. I expect
125\% a year, unlevered.'' Give three independent reasons why the expectation
is wrong.
\end{exercise}

\section{Problem: Reconstitution Day}

\begin{problem}[{Weekend problem --- one small stock on the last Friday of June}]\label{pb:m1:index-construction-and-rebalancing:1}
A small-capitalisation index has a float-adjusted capitalisation of \$3
trillion and trackers with \$200 billion. A large-capitalisation index of the
same family has \$50 trillion and trackers with \$4 trillion. Stock X (price
\$24, float-adjusted capitalisation \$1.2 billion, average volume 0.9 million
shares a day, daily volatility 3\%) is to join the small index. Stock Y (price
\$60, float \$9 billion, average volume 2.5 million shares) migrates from the
small index to the large one.

\medskip\noindent\textbf{Part I --- The demand.}
\begin{enumerate}
\item What fraction of each member's float do the small index's trackers
hold? And the large index's?
\item Give the trackers' demand for X in dollars and in shares.
\item Express it in days of average volume.
\item For Y, give the small trackers' sale, the large trackers' purchase and
the net.
\item Express the sale and the net in days of Y's volume. Which number
matters for the closing price?
\end{enumerate}

\medskip\noindent\textbf{Part II --- The close.}
\begin{enumerate}[resume]
\item Trackers send 60\% of their demand for X to the closing auction. X's
closing auction normally trades 8\% of a day's volume. Compare.
\item Estimate the impact of the whole demand for X with the square-root
rule, $\kappa = 0.7$, as if it arrived in one day.
\item Same if spread evenly over twenty days.
\item Who sells 2 million shares of X in the auction, and where did they get
them?
\end{enumerate}

\medskip\noindent\textbf{Part III --- The anticipation trade.}
Three weeks before the rank day a firm believes X will be added with
probability 0.8. If added, it expects a run-up of 3\% to the effective close;
if not, a fall of 2\% as other anticipators leave. A round trip costs 40
basis points. It buys \$5 million.
\begin{enumerate}[resume]
\item Give the expected return and the expected profit.
\item Give the standard deviation of the profit.
\item The firm holds 25 such positions. With independent outcomes, give the
expected profit, its standard deviation and their ratio.
\item Same with a pairwise correlation of 0.3 between the positions' returns.
\item At what probability of addition does the trade break even?
\item Name two sources of that correlation.
\end{enumerate}

\medskip\noindent\textbf{Part IV --- Judgement.}
\begin{enumerate}[resume]
\item The provider moves from one reconstitution a year to two. What happens
to the size of each event and to the anticipators' business?
\item Why do the trackers accept to pay the impact at the close instead of
buying early?
\item An index fund's manager \emph{does} buy early. Who bears the risk, and
how would its investors find out?
\item State the \emph{named result}: the trackers' demand for X as a multiple
of its average daily volume.
\item In one sentence: what is the economic function of the firms that
anticipate index changes?
\end{enumerate}
\end{problem}

\section{Interview questions}

\begin{interviewq}[$\star$ \iqroles{trader, researcher}]\label{iq:m1:index-construction-and-rebalancing:1}
A stock is added to a major index. What happens to its price, and when?
\end{interviewq}

\begin{interviewq}[$\star$ \iqroles{researcher, developer}]\label{iq:m1:index-construction-and-rebalancing:2}
Why are indices float-adjusted, and what breaks if they are not?
\end{interviewq}

\begin{interviewq}[$\star\star$ \iqroles{researcher, trader}]\label{iq:m1:index-construction-and-rebalancing:3}
Trackers hold 8\% of an index. A stock with a \$10 billion float is added.
How much must they buy? What else do you need to know to say whether that is
a lot?
\end{interviewq}

\begin{interviewq}[$\star\star$ \iqroles{researcher}]\label{iq:m1:index-construction-and-rebalancing:4}
How would you predict next June's additions to a rules-based small-cap index,
and where would your errors come from?
\end{interviewq}

\begin{interviewq}[$\star\star$ \iqroles{researcher, trader}]\label{iq:m1:index-construction-and-rebalancing:5}
Passive assets have grown enormously, yet the measured \omterm{def:m1:index-construction-and-rebalancing:effect}{index effect} has
shrunk. Reconcile.
\end{interviewq}

\begin{interviewq}[$\star\star\star$ \iqroles{trader, researcher}]\label{iq:m1:index-construction-and-rebalancing:6}
You run an index-rebalance book. Describe its main risks and how you would
size positions.
\end{interviewq}
